% Independently audited by root and the F/topos agent; integrated as AG.39.
\subsection{A failure of full prime-to-\texorpdfstring{\(p\)}{p} configuration exactness}

\begin{claim}{AG.39}{Proved: the full prime-to-\(p\) configuration fiber does not inject}
\label{clm:AGconfigurationodd}
Let \(C\) be a smooth proper connected complex curve of genus
\(g\geq2\), let \(P\in C\), and put
\[
 U=\{(x,z)\in C\times C:x\ne z\},\qquad F=C-\{P\}.
\]
Use the second projection \(U\to C\) to identify \(F\) with a
fiber. For every prime \(p\), write \(\Gamma^{(p')}\) for the
inverse limit of the finite quotients of a group \(\Gamma\)
whose orders are prime to \(p\). Then
\[
 \pi_1(F)^{(p')}\longrightarrow
 \pi_1(U)^{(p')}                                     \tag{AG.39.1}
\]
is not injective. The same assertion holds for the geometric
prime-to-\(p\) etale fundamental groups of every smooth proper
curve of genus at least two in characteristic \(p\).
In particular it holds for the smooth projective model of
\[
                    C_0/\mathbb F_2:\quad y^2+y=x^5,
             \qquad P=\infty.                       \tag{AG.39.2}
\]

This refutes the full prime-to-\(p\) configuration exactness
condition proposed as a sufficient route in
\cite[Remark 3.13]{AR_ST2009}. It does not supply two coincident
point sections and does not refute A1. The proof specifies a finite
curve and finite branched-cover construction; it does not enumerate
a particular nonidentity profinite kernel element or an individual
finite monodromy matrix witnessing the failure.
\end{claim}

\begin{proof}[Independent proof for \(p=2\)]
For this proof write \((-)^{\mathrm{odd}}=(-)^{(2')}\).
The extension to every prime is proved below by a separate
congruence argument.
\emph{An odd-degree branched family.}
Choose a connected cyclic etale cover
\(a:C'\to C\) of degree three. Put \(m=3^{2g}\), and let
\(b:B\to C\) be the connected etale cover corresponding to
\[
 \ker\bigl(\pi_1(C)\longrightarrow H_1(C,\mathbb F_3)\bigr).
                                                               \tag{AG.39.3}
\]
It is Galois of degree \(m\), factors through \(C'\), and
\(b^*:H^1(C,\mathbb F_3)\to H^1(B,\mathbb F_3)\) is zero:
each homomorphism \(\pi_1(C)\to\mathbb F_3\) restricts to
zero on the subgroup in (AG.39.3).

On \(Y=C'\times B\) consider the effective divisor
\[
          D=(a\times b)^{-1}(\Delta_C).
                                                               \tag{AG.39.4}
\]
It consists of three disjoint graphs \(B\to C'\). Its first
Chern class is divisible by three in \(H^2(Y,\mathbb Z)\).
Indeed, in the Kunneth decomposition modulo three, the two
degree-two summands of the diagonal class pull back by maps of
degrees \(3\) and \(m\), and its cross summand vanishes because
\(b^*H^1(C,\mathbb F_3)=0\). The integral cohomology of a
product of compact oriented surfaces is torsion-free, so vanishing
modulo three is equivalent to integral divisibility.
The class \(c_1(D)/3\) is integral of type \((1,1)\); the
Lefschetz \((1,1)\) theorem gives a line bundle with that Chern
class. Its cube differs from \(\mathcal O_Y(D)\) by an element
of \(\operatorname{Pic}^0(Y)\). Surjectivity of multiplication
by three on this complex abelian variety therefore gives
\[
                         L^{\otimes3}\simeq\mathcal O_Y(D).
                                                               \tag{AG.39.5}
\]

Take the cyclic cover \(S\to Y\) obtained by adjoining a cube
root of the canonical section of \(L^{\otimes3}\). It is a
connected smooth projective surface. The composite
\(f:S\to B\) is smooth and proper with connected fibers:
locally along any graph the cover has the form \(w^3=t\), where
\((t,z)\) are coordinates with \(z\) on \(B\); these become
the coordinates \((w,z)\) on \(S\). Each fiber is a cyclic
triple cover of \(C'\), branched at the three distinct points
\(a^{-1}(b(z))\), each with ramification index three. Consequently
\[
 g(C')=3g-2,\quad g(B)=m(g-1)+1,\quad
                  h:=g(S_z)=9g-5.                   \tag{AG.39.6}
\]
This is the usual cyclic construction of a Kodaira fibration;
compare \cite[Section 1.3]{AG_Rollenske2010}. All the particular
divisibility and smoothness assertions needed here were proved
above.

The family is nonisotrivial. To see this without any genericity
assumption, suppose every fiber were isomorphic to a fixed curve
\(T\). There are only finitely many degree-three morphisms
\(T\to C'\): their Hom scheme is of finite type, and the tangent
space at a morphism \(q\) is
\(H^0(T,q^*T_{C'})=0\), since this line bundle has negative
degree. Thus there would be only finitely many possible branch
sets in \(C'\). The sets \(a^{-1}(b(z))\), however, vary with
the surjective map \(b:B\to C\). This is a contradiction.

\emph{The required comparison of completions.}
Let \(S^\circ\) be the complement of the ramification divisor
in \(S\). There is a tower of finite etale covers
\[
 S^\circ\longrightarrow Y-D
 \longrightarrow (C\times B)-\operatorname{graph}(b)
 \longrightarrow U.                                 \tag{AG.39.7}
\]
The first two arrows are cyclic of degree three, and the last is
the Galois base change of (AG.39.3). The Galois closure of the
composite has a three-group as its group: the Galois closure of a
tower of covers with three-group Galois groups embeds in the
corresponding iterated wreath product. In a fiber of \(U\to C\),
the cover \(S_z^\circ\to F\) has degree nine and likewise has
three-group Galois closure.

Here is the elementary topology-of-completion fact being used.
If \(H\leq E\) contains a normal subgroup of finite odd index,
then \(H^{\mathrm{odd}}\) is the open subgroup given by the
closure of \(H\) in \(E^{\mathrm{odd}}\). For any normal
odd-index subgroup \(N\) of \(H\), first intersect \(N\)
with that fixed normal subgroup of \(E\), and then take its
\(E\)-core. The resulting quotient is an extension of an odd
group by a subgroup of a finite product of odd groups. It has
odd order and its restriction to \(H\) refines the given
quotient. This proves fullness of the induced pro-odd topology.

Suppose now that (AG.39.1) were injective. Apply the preceding
fact to both the total cover and its degree-nine fiber cover.
The resulting homomorphism
\[
 \pi_1(S_z^\circ)^{\mathrm{odd}}\longrightarrow
                 \pi_1(S^\circ)^{\mathrm{odd}}        \tag{AG.39.8}
\]
is injective. The discrete fiber group is normal in the discrete
total group and the latter surjects onto \(\pi_1(B)\):
\(S^\circ\to B\) is a locally trivial punctured-surface
fibration with connected fiber. Thus its compact pro-odd closure
is normal in the completed total group. These assertions also
follow from the Fadell--Neuwirth sequence for \(U\to C\) and
the pullback covers; see \cite[Section 2.1]{AG_Bellingeri2004}.

For every odd prime \(\ell\), compactification of the punctured
fiber gives a quotient
\[
 \pi_1(S_z^\circ)^{\mathrm{odd}}\longrightarrow
                          H_1(S_z,\mathbb F_\ell).
                                                               \tag{AG.39.9}
\]
Its kernel is stable under the discrete total group, because the
family preserves the peripheral loops, and therefore under its
pro-odd closure. The conjugation action on this finite quotient
is continuous. For example this follows by testing the images
of a finite generating set of the pro-odd fiber group. Hence its
image has odd order. Inner conjugation by the fiber group acts
trivially on (AG.39.9). Since the total discrete group surjects
onto \(\pi_1(B)\), the image just obtained is precisely the
usual monodromy image on \(H_1(S_z,\mathbb F_\ell)\).
In particular, if
\[
 \Gamma=\operatorname{im}\bigl(\pi_1(B)\to
                    \operatorname{Sp}_{2h}(\mathbb Z)\bigr),
\]
then the reduction of \(\Gamma\) modulo every odd prime has
odd order. No assertion that completion is left exact for an
arbitrary surface fibration has been used.

\emph{The monodromy contradiction, in characteristic zero.}
The Feit--Thompson theorem makes each of those finite reductions
solvable \cite{AC_FeitThompson1963}. The uniform solvable-length
theorem for linear groups supplies an integer \(d(2h)\),
depending only on the dimension and not on the field, that
bounds their derived lengths; see
\cite[Section 1, printed p. 59]{AG_Newman1972}.
Every element of \(\Gamma^{(d(2h))}\) therefore becomes the
identity modulo every odd prime. An integral matrix with this
property is the identity. Thus \(\Gamma\) is solvable.

On the other hand, Deligne's theorem states that the identity
component of the Zariski closure of the monodromy of a smooth
proper complex algebraic family is semisimple
\cite[Corollary 4.2.9(a)]{AG_Deligne1971}.
The Zariski closure of a solvable group is solvable, so that
identity component must be trivial. Hence \(\Gamma\) is finite.
After a finite etale cover of \(B\) the integral local system
is constant. Its period map is then a globally defined
holomorphic map from a compact curve into Siegel upper half
space. In its bounded realization all coordinate functions are
constant. The principally polarized Jacobian is constant, and
Torelli makes the fibers of \(f\) isomorphic. This contradicts
the nonisotriviality already proved and proves (AG.39.1).

\emph{Passage to characteristic two.}
This last Hodge argument is used only over \(\mathbb C\).
Let \((C,P)\) be a pointed smooth proper curve over an
algebraically closed field of characteristic two. Choose a
finitely generated subfield over \(\mathbb F_2\) over which
\((C,P)\) is defined, and replace it by its algebraic closure
\(k_0\) inside the given field. Prime-to-two fundamental groups
of these smooth proper pairs are invariant under algebraically
closed extension, so it suffices to work over \(k_0\).
A pointed smooth proper lifting over \(W(k_0)\) exists: the
obstruction group for deforming a pointed smooth curve is
\(H^2(C,T_C(-P))=0\), and a relatively ample line bundle
algebraizes the formal lifting. The countable field \(k_0\)
has a Witt fraction field of cardinality at most that of
\(\mathbb C\); an algebraic closure of this characteristic-zero
field therefore embeds abstractly into \(\mathbb C\). This
gives the complex geometric generic fiber used below.

For this lift, both
\((\mathcal C,\mathcal P)\) and
\((\mathcal C\times\mathcal C,\Delta_{\mathcal C})\)
are smooth proper pairs with relative normal-crossing boundary.
SGA~1 XIII, Corollary~2.9 and the specialization construction in
Section~2.10 identify their maximal prime-to-two fundamental
groups on the geometric generic and special fibers
\cite{AR_SGA1}. In more detail, Corollary~2.9 makes the nonabelian
\(R^1\) for every finite group of odd order locally constant;
its base-change statement and Section~2.10 give these
specialization isomorphisms. Functoriality for
\(x\mapsto(x,\mathcal P)\) identifies the fiber-inclusion maps
as well. Injectivity in characteristic two would therefore imply
injectivity for the complex curve, which was disproved above.

For the completely specified example (AG.39.2), a lift requires
no deformation argument. The equation \(y^2+y=x^5\) over
\(\mathbb Z_2\), with the chart at infinity
\[
                         v^2+u^3v=u,
\]
gives the required smooth genus-two model and the section
\(u=v=0\). On the special fiber the finite chart is smooth
because its derivative with respect to \(y\) is one; the
infinity chart is smooth at that section because its derivative
with respect to \(u\) is a unit. Its generic fiber embeds into
\(\mathbb C\), so the same functorial specialization diagram
applies.
\end{proof}

\paragraph{Finite numerical checks and scope.}
For genus two, (AG.39.3)--(AG.39.7) have
\[
 m=81,\qquad g(C')=4,\qquad g(B)=82,
 \qquad g(S_z)=13,
\]
and the punctured fiber has three punctures and first Betti
number \(28\); its compactification has first Betti number
\(26\). The retained exact-rational program
\path{configuration_family_invariants.py} also verifies
\[
 D^2=-486,\quad K_YD=972,\quad K_Y^2=1944,
 \quad e(S)=3888,\quad K_S^2=9072,\quad
 \operatorname{sign}(S)=432.
                                                               \tag{AG.39.10}
\]
The signature is an independent numerical consistency check,
not a substitute for the monodromy argument. The new obstruction
concerns completing the \emph{whole} configuration group to
pro-odd. It does not contradict the exact sequence with full
tame base and only the fiber replaced by its prime-to-two
quotient, nor the exact single-prime pro-\(\ell\) sequences.

\paragraph{A cohomological obstruction to naive cuspidal devissage.}
There is a concrete loss of information when one tries to
iterate maximal cuspidally pro-\(\ell\) reconstruction through
ramified covers. Work over an algebraically closed field of
characteristic two, or over \(\mathbb C\), and take the
pro-odd groups throughout. Let \(C\) have genus \(g\geq2\),
let \(P\in C\), and write
\[
 F=\pi_1(C-P)^{\mathrm{odd}},\quad
 G=\pi_1(C)^{\mathrm{odd}},\quad K=\ker(F\to G),
 \quad N=\ker(K\to K^{(3)}).
\]
Here \(K^{(3)}\) denotes the maximal pro-three quotient,
so \(F/N\) is the maximal cuspidally pro-three quotient over
\(G\). Choose a cyclic etale cubic \(a:C'\to C\), and a
cyclic cubic \(D\to C'\) branched exactly at the three
points of \(a^{-1}(P)\), with ramification index three at
each. Such a cover exists because the degree-three branch
divisor has a cube root in \(\operatorname{Pic}(C')\):
choose a line bundle of degree one and correct its cube
by divisibility of \(\operatorname{Pic}^0(C')\) over the
algebraically closed field.
It has genus \(h=9g-5\).

Put \(D^\circ=D-\pi^{-1}(P)\). Its degree-nine etale cover
of \(C-P\) is visible in \(F/N\), since its Galois closure
has three-group Galois group. Let \(H\leq F/N\) be its
open subgroup, and let \(\overline H\) be the quotient of
\(H\) by the closed normal subgroup generated by its three
puncture inertia groups. Then, for every odd prime
\(\ell\ne3\),
\[
 \dim_{\mathbb F_\ell}H^1(\overline H,\mathbb F_\ell)
       =2g(C')=6g-4
       <18g-10
       =\dim_{\mathbb F_\ell}H^1(D,\mathbb F_\ell).
                                                               \tag{AG.39.11}
\]
Thus \(\overline H\) is a proper quotient of the full
group \(\pi_1(D)^{\mathrm{odd}}\), already detected on
finite abelian quotients. For \(g=2\), the two dimensions
are \(8\) and \(26\).

Here is a direct proof of the comparison, including its
surjectivity assertion. Let \(F'\leq F\) correspond to
\(C'-a^{-1}(P)\), and let
\(\chi:F'\twoheadrightarrow\mathbb Z/3\) define
\(D^\circ\). The group \(K\) is contained in \(F'\),
and \(\chi|_K\) is surjective, because a branch inertia
generator lies in \(K\) and maps nontrivially. If
\(I=\ker\chi\), this gives \(F'=IK\). Also \(N\subset I\),
and hence
\[
 1\longrightarrow (I\cap K)/N\longrightarrow
 H=I/N\longrightarrow \pi_1(C')^{\mathrm{odd}}
 \longrightarrow1.                                  \tag{AG.39.12}
\]
The left-hand term is a closed subgroup of the pro-three
group \(K/N\). Every homomorphism from \(H\) to
\(\mathbb F_\ell\) therefore factors uniquely through
the last group when \(\ell\ne3\). The puncture inertia
groups in \(H\) lie in the left-hand term, so killing them
does not change these homomorphisms. This proves the first
equality in (AG.39.11); the second is the usual surface
cohomology dimension and the computed genus of \(D\).
The natural map from \(\pi_1(D)^{\mathrm{odd}}\) to
\(\overline H\) is surjective by the construction.

This is an obstruction to the specific procedure of taking
an open subgroup of a cuspidally pro-three quotient and then
killing its puncture inertia. It is not a theorem that all
cuspidal quotients together determine insufficient information
by every possible method. In a solvable induction, the next
cuspidally pro-\(\ell\) theorem needs the \emph{full} proper
prime-to-two group of the intermediate ramified compactification.
Equation (AG.39.11) shows exactly why the preceding formal
operations do not provide that input. That ramified
compactification is also absent from the original proper
unramified tower.

% Extension of AG.39; intended for incorporation into that statement/proof,
% not a second independently numbered configuration claim.

\paragraph{The congruence lemma needed for every excluded prime.}
Let \(\Gamma\subset\operatorname{GL}_N(\mathbb Z[1/M])\) be
finitely generated. Suppose that the identity component \(G\) of its
Zariski closure over \(\mathbb Q\) is a nontrivial semisimple group.
For every prime \(p\) and every integer \(a\geq1\), there are
infinitely many rational primes \(\ell\ne p\), with \(\ell\nmid M\),
such that
\[
       p^a\mid
       \#\operatorname{im}\bigl(\Gamma\longrightarrow
                    \operatorname{GL}_N(\mathbb F_\ell)\bigr).
                                                        \tag{AG.39.13}
\]

\begin{proof}
The subgroup \(\Gamma_0=\Gamma\cap G(\mathbb Q)\) has finite index,
is finitely generated, and is Zariski dense in \(G\). Let
\[
     1\longrightarrow Z\longrightarrow \widetilde G
          \xrightarrow{\pi}G\longrightarrow1
\]
be its simply connected central covering, and put \(c=\#Z\).
The connecting map
\[
               \delta:G(\mathbb Q)\longrightarrow H^1(\mathbb Q,Z)
\]
is a group homomorphism. Indeed, for lifts of two rational elements,
the two Galois cocycles multiply because their values lie in the
central subgroup \(Z\). The abelian group \(H^1(\mathbb Q,Z)\)
is killed by \(c\). Therefore the image of the finitely generated
group \(\Gamma_0\) under \(\delta\) is finite. Its kernel
\(\Gamma_1\) has finite index in \(\Gamma_0\), and every element
of \(\Gamma_1\) has a rational lift through \(\pi\).

Choose a finite generating set of \(\Gamma_1\), choose rational
lifts of its generators, and let \(\widetilde\Gamma\) be the
group generated by those lifts. Its image is exactly \(\Gamma_1\).
Its Zariski closure is all of \(\widetilde G\): the image of that
closure under the finite morphism \(\pi\) is closed and contains
the dense subgroup \(\Gamma_1\); equality of dimensions, and
connectedness of \(\widetilde G\), finish the assertion.

Choose a faithful rational representation of \(\widetilde G\).
After enlarging a finite set of primes \(S\), all chosen lifts
are integral over \(\mathbb Z[S^{-1}]\), and \(\pi\), the
original faithful representation of \(G\), and their central
kernel spread out over this ring. We may require that these
models are smooth semisimple group schemes, that the morphism
is a central isogeny of degree \(c\), and that the embedding
of \(G\) into the original matrix group is a closed immersion.
Thus, outside \(S\), reducing \(\pi(\widetilde\gamma)\) gives
the same matrix as reducing the original element of \(\Gamma_1\).
The kernel of this comparison on finite-field points has at
most \(c\) elements.

Nori's strong-approximation theorem, in precisely its
connected and simply connected form
\cite[Theorem 5.4, printed p. 274]{AC_Nori1987}, says that the
congruence closure of \(\widetilde\Gamma\) is open in
\(\widetilde G(\prod_{\ell\notin S}\mathbb Z_\ell)\).
Consequently, after excluding finitely many more primes,
\[
       \widetilde\Gamma\bmod\ell
                      =\widetilde G(\mathbb F_\ell).
\]
Here smoothness gives surjectivity of reduction from
\(\mathbb Z_\ell\)-points, and an open subgroup of a product
contains the full product outside finitely many factors.

Choose a finite Galois number field \(K\) which splits
\(\widetilde G\) and contains the roots of unity of order
\(p^{a+v_p(c)}\). Exclude also the finitely many rational
primes needed to spread the chosen splitting isomorphism
over the ring of integers of \(K\). Infinitely many rational primes split
completely in \(K\). One may use the usual splitting-prime
theorem, or the following elementary sufficient form. If
\(K=\mathbb Q(\alpha)\) with \(\alpha\) integral and minimal
polynomial \(f\), infinitely many primes divide nonzero values
\(f(n)\). Indeed, if all such prime divisors belonged to a
finite set with product \(R\), choose \(b\) with
\(d=f(b)\ne0\). The integer polynomial
\(f(b+dRt)/d\) is congruent to one modulo \(R\) and has
nonconstant absolute value, a contradiction for sufficiently
large integers \(t\). Away from the discriminant and the
index of \(\mathbb Z[\alpha]\), a root of \(f\) modulo a
prime gives a residue-degree-one prime of \(K\). Since
\(K/\mathbb Q\) is Galois, that rational prime splits completely.

Discard the finite exceptional set, and choose any such
\(\ell\ne p\). Then \(\widetilde G_{\mathbb F_\ell}\)
is a nontrivial split semisimple group and
\(p^{a+v_p(c)}\mid\ell-1\). Its split maximal torus has
positive rank, so
\[
       v_p\bigl(\#\widetilde G(\mathbb F_\ell)\bigr)
                          \geq a+v_p(c).
\]
Its image in the original \(\operatorname{GL}_N(\mathbb F_\ell)\)
loses at most \(v_p(c)\) from this valuation, and that image is
the reduction of \(\Gamma_1\). It is a subgroup of the
reduction of \(\Gamma\). This proves (AG.39.13).
\end{proof}

\paragraph{Replacement of the pro-odd restriction in AG.39.}
Fix an arbitrary prime \(p\), and write \((-)^{(p')}\) for
completion with respect to all finite groups of order prime
to \(p\). For every smooth proper complex curve \(C\) of
genus \(g\geq2\), every \(P\in C\), and
\(U=(C\times C)-\Delta_C\), the map
\[
       \pi_1(C-\{P\})^{(p')}
                 \longrightarrow \pi_1(U)^{(p')}
                                                        \tag{AG.39.14}
\]
is not injective. The corresponding geometric prime-to-\(p\)
map is likewise noninjective for every smooth proper curve of
genus at least two over an algebraically closed field of
characteristic \(p\).

\begin{proof}
In the explicit family construction of AG.39 choose a prime
\(r\ne p\): one may take \(r=2\) for odd \(p\), and
\(r=3\) for \(p=2\). Replace its cyclic triple cover by a
cyclic degree-\(r\) etale cover \(a:C'\to C\), and take
\(b:B\to C\) to be the full mod-\(r\) homology cover,
of degree \(r^{2g}\). This kills the cross term of the
diagonal class modulo \(r\), and both remaining Kunneth
terms vanish modulo \(r\). The same integral-cohomology and
Picard-divisibility argument therefore gives
\[
       L^{\otimes r}\simeq
       \mathcal O_{C'\times B}\bigl((a\times b)^{-1}\Delta_C\bigr).
\]
The resulting cyclic cover is a smooth proper family
\(S\to B\). Each fiber is a degree-\(r\) cover of \(C'\)
branched at the \(r\) distinct points \(a^{-1}(b(z))\), and
\[
   g(C')=1+r(g-1),\qquad g(B)=1+r^{2g}(g-1),\qquad
   g(S_z)=1+r^2(g-1)+\frac{r(r-1)}2.
\]
It is nonisotrivial by the same finite-Hom argument: a fixed
source curve has only finitely many degree-\(r\) maps to the
fixed hyperbolic curve \(C'\), while their required branch
sets vary through the fibers of the surjection \(a\).

Every cover in the punctured total-space tower has an
\(r\)-group Galois group, and its composite has an
\(r\)-group Galois closure by the wreath-product argument.
The punctured fiber has the same property. An open subgroup
containing a normal subgroup of finite prime-to-\(p\) index
has its full pro-prime-to-\(p\) topology induced from the
ambient group: taking normal cores produces an extension
of a prime-to-\(p\) group by a subgroup of a product of such
groups. This uses no solvability assumption.

Consequently, if (AG.39.14) were injective, the completed
punctured-fiber group would be a closed normal subgroup of
the completed total group of \(S^\circ\to B\). For every
\(\ell\ne p\), its compactified homology
\(H_1(S_z,\mathbb F_\ell)\) is a finite quotient with
monodromy-invariant kernel. Conjugation is continuous, the
fiber acts trivially on this abelian quotient, and the
discrete total group surjects onto \(\pi_1(B)\). Therefore
every reduction modulo \(\ell\ne p\) of the integral
monodromy group
\[
       \Gamma\subset\operatorname{Sp}_{2g(S_z)}(\mathbb Z)
\]
would have order prime to \(p\). These are reductions of
the actual integral homology local system, rather than
abstract finite quotients of an unrelated group.

Deligne's semisimplicity theorem used in AG.39 makes the
connected Zariski closure of \(\Gamma\) semisimple. It is
nontrivial: a trivial identity component would make
\(\Gamma\) finite, so after a finite etale base change the
period map would be a holomorphic map from a compact curve
to the bounded realization of Siegel space. It would be
constant; Torelli would make the family isotrivial. The
congruence lemma, with \(a=1\), now contradicts the asserted
prime-to-\(p\) orders of all the reductions. This proves
(AG.39.14) over \(\mathbb C\).

For characteristic \(p\), lift the pointed smooth curve
as in AG.39 and apply the prime-to-\(p\) specialization
isomorphisms for the two smooth proper pairs
\((\mathcal C,\mathcal P)\) and
\((\mathcal C\times\mathcal C,\Delta_{\mathcal C})\).
SGA~1 XIII, Corollary~2.9 and Section~2.10 apply because
all the finite groups in question have order invertible
on the base. Functoriality identifies the two fiber maps.
Thus injectivity of the special-fiber map would imply
injectivity of its complex counterpart, which has just
been contradicted. The Hodge and period-map argument has
been used only in characteristic zero.
\end{proof}

The extension eliminates the full configuration exactness
shortcut for every excluded prime. It does not produce a
collision, and by itself does not decide A1.
The later AG40--AG42 settle A1; universal relative-witness
existence remains a separate question. The characteristic-two proof in AG.39
via Feit--Thompson and bounded solvable length remains an
independent alternative to the strong-approximation step.

\paragraph{What the same-prime quotient does retain.}
The preceding cross-prime rank loss must not be confused
with a loss of the entire zeta function. There is a general
elementary complement. Let \(F\) be a profinite group,
let \(K\triangleleft F\) be closed, fix a prime \(r\),
and put \(N=\ker(K\to K^{(r)})\). For every open
subgroup \(I\leq F\) containing \(N\), the natural map
induces an isomorphism
\[
                         I^{(r)}\simeq (I/N)^{(r)}.
                                                        \tag{AG.39.15}
\]
Here the superscript denotes the maximal pro-\(r\)
quotient. In particular, the assertion is not restricted
to open subgroups having an \(r\)-group Galois core.

To prove it, take an open normal subgroup \(J\triangleleft I\)
such that \(I/J\) is an \(r\)-group, and set
\(I_0=\operatorname{core}_F(I)\).
The group \(J\cap I_0\) is normal in \(I_0\), and
\(I_0/(J\cap I_0)\) embeds in \(I/J\).
Let \(M=\operatorname{core}_F(J\cap I_0)\).
Then \(I_0/M\) embeds in a finite product of
\(r\)-groups, so it is an \(r\)-group. Since
\(N\triangleleft F\) and \(N\subset I\), we have
\(N\subset I_0\); consequently the image of \(K\)
in \(F/I_0\) is an \(r\)-group. Its image in
\(F/M\) is an extension of an \(r\)-group by a
subgroup of \(I_0/M\), and is therefore an
\(r\)-group as well. The defining property of \(N\)
now gives \(N\subset M\subset J\). Thus every
finite \(r\)-quotient of \(I\) kills \(N\),
which proves (AG.39.15).

Apply this to the geometric punctured-curve group and
its cuspidal kernel. For every finite cover visible in
the maximal cuspidally pro-\(r\) quotient, quotienting
by the puncture inertia and then taking the maximal
pro-\(r\) quotient recovers the \emph{full} proper
pro-\(r\) group of its smooth compactification.
Indeed, formation of a maximal pro-\(r\) quotient
commutes with quotienting by the closed normal subgroup
generated by the specified inertia groups. This applies
in particular to the ramified cover \(D\) in
(AG.39.11), even though its other-prime cohomology has
been lost.

If the curves and cover are defined over a finite field
of characteristic different from \(r\), this comparison
is compatible with the arithmetic Frobenius action.
It therefore retains \(H^1(D_{\bar k},\mathbb Z_r)\)
with Frobenius, and hence the exact integral Weil
polynomial and zeta function of \(D\). Equality of the
\(r\)-adic characteristic polynomials is equality of
their integer polynomials, since
\(\mathbb Z[T]\hookrightarrow\mathbb Z_r[T]\).
Thus the obstruction in (AG.39.11) concerns the full
proper prime-to-\(p\) group needed by a subsequent
different-prime reconstruction theorem; it does not
prevent transfer of these zeta functions. No inference
from that transfer to point separation is made here.
In particular, the Weil polynomial does not by itself
specify the integral \(p\)-primary rational torsion
exponent, or supply R1-compatible point lifts on these
ramified covers.

